\documentclass[10pt,twocolumn]{article}
\usepackage[a4paper,margin=1.9cm,columnsep=0.7cm]{geometry}
\usepackage{mathptmx}
\usepackage[T1]{fontenc}
\usepackage{microtype}
\usepackage{graphicx}
\usepackage{booktabs}
\usepackage{caption}
\usepackage{authblk}
\usepackage[hidelinks]{hyperref}
\usepackage{xcolor}

\captionsetup{font=small,labelfont=bf,skip=5pt}
\raggedbottom
\setlength{\parindent}{1.1em}
\definecolor{mono}{HTML}{2F6F9F}
\definecolor{dia}{HTML}{C25A3C}
\definecolor{m3gray}{HTML}{8A8A8A}

\title{\vspace{-1.2em}\bfseries Monologue or Dialogue?\\[2pt]
\large Measuring Conversational Structure in the Reasoning Traces of Large Language Models}
\author[ ]{Oskar Engl}
\affil[ ]{\small Johannes Kepler University Linz, Austria \quad\textit{oskar.engl@ik.me}}
\date{\small Working version, July 2026}

\begin{document}
\twocolumn[
\begin{@twocolumnfalse}
\maketitle
\begin{abstract}
\noindent
A reasoning model such as DeepSeek-R1 produces a long ``thinking'' trace before it gives an answer.
That trace can be read as one mind working step by step, or as several voices exchanging proposals
and objections. This paper asks which reading the text itself supports once it is measured. No
reasoning trace comes labelled with its true voice structure, so we build a small grid of corpora
whose structure is known: human and AI, monologue and dialogue. We check that each of three
text-analysis methods can separate monologue from dialogue on these known corpora, and then apply
them to the traces. Two of the three place the traces on the dialogue side of the scale, and the
third, which reads the coarsest signal, places them just onto the monologue side by a small margin.
The three methods are not equally well validated, so we state how much weight each one can bear and
flag the one that depends on an automatic labeller. Taken together, the traces show several surface
marks of dialogue, most clearly in how the writing style shifts from sentence to sentence and in how
often a question is followed by an answer.
\medskip
\end{abstract}
\end{@twocolumnfalse}
]

\begin{figure*}[t]
\centering
\includegraphics[width=0.96\textwidth]{fin_fig_grid.pdf}
\caption{\textbf{Study design.} Four corpora of known voice structure (left) feed a validity check
that each method must pass: separating monologue from dialogue on all four known cells. Methods that
pass are applied to the DeepSeek-R1 traces (\texttt{ai-reason}) and place them on a
monologue--dialogue scale anchored by the two same-task AI cells; the human cells serve only to
confirm that a method works. The check shown here is the monologue/dialogue separation, which all
three methods pass. Method~3 faces a further check on the reliability of its labels, described in
Section~4.}
\label{fig:grid}
\end{figure*}

\section{The question}
Reasoning-oriented training has become standard in frontier models. Given a problem, a model such as
DeepSeek-R1 \cite{guo2025} first produces a long stretch of intermediate text, often called a
reasoning trace \cite{kambhampati2025}, and only then gives its final answer. These traces have a
notable feature: each is written by a single model in one pass, yet a trace often contains questions
and their answers, objections and concessions, and sudden reversals such as ``Wait, that is not
right.'' In writing by people, features like these can mark a real exchange between participants,
though they can also mark a single inner voice. This raises a question that is easy to pose and hard
to settle: does a reasoning trace have the textual shape of a monologue or of a dialogue?

The question also bears on current theory. Kim et al.~\cite{kim2026} argue that reasoning-trained
models improve in part because they simulate a ``society of thought,'' a set of internal perspectives
that propose, criticise, and reconcile. Older ideas in cognitive science point the same way, from
reasoning as a social and argumentative activity \cite{mercier2017} to the mind as a set of
interacting agents \cite{minsky1986}. Our premise is narrower and testable: if a trace is produced by
several voices taking turns, that should leave marks in the writing. Writers differ in small and
largely involuntary habits, such as how they use function words, how their sentences flow, and which
conversational moves they make, and authorship analysis and conversation analysis use these habits to
tell writers and turns apart. We therefore treat the question in operational terms: a monologue
should read as the work of one author, and a dialogue as several voices exchanging turns, and we ask
which of these the traces resemble.

\section{The design: a calibration grid}
No collection of reasoning traces comes labelled with its true voice structure, which is what makes
the question hard to answer directly. We address this with the grid in Figure~\ref{fig:grid}. We
assemble four corpora whose structure is known and vary two properties across them: who produced the
text (a human or an AI) and whether it is a monologue or a dialogue. The four cells are human
monologue (think-aloud protocols, in which a person speaks their thoughts aloud while solving a task
\cite{wurgaft2025}), human dialogue (group deliberation from DeliData \cite{karadzhov2023}), AI
monologue (a single Qwen-2.5-7B model solving a puzzle on its own), and AI dialogue (four separate
copies of Qwen-2.5-7B, each given a role, taking turns in a debate). The fifth cell, the DeepSeek-R1
traces, is the one whose structure we want to determine.

The grid functions as a validity check, or gate. A method may report on the traces only after it has
shown that it can separate monologue from dialogue on all four known corpora. The rationale is
straightforward: a method that cannot tell the voices apart in a genuine four-way debate is in no
position to rule them out in a trace, and a method that reports several voices in one person's
think-aloud protocol is in no position to affirm them there.

One design choice keeps the comparison clean: we hold the task fixed where possible. The AI
monologue, AI dialogue, and DeepSeek-R1 cells all solve the same puzzle, Game of 24 \cite{yao2023}, a
small arithmetic puzzle in which four given numbers must be combined with the four basic operations to
reach 24, and the human monologue cell uses the same puzzle. A method that reads a monologue as a
single voice therefore cannot be responding merely to the subject matter. The human dialogue cell is
the one exception, since it uses the Wason card task, another compact logic puzzle, because no
group-deliberation corpus exists for Game of 24; its role is only to confirm that a method can detect
human dialogue at all, so it never anchors the placement. We also removed speaker names, role titles,
and letter labels from both dialogue corpora before analysis, so that a method cannot detect dialogue
by identifying who is speaking instead of reading the structure of the exchange.

\section{Three methods, three signals}
The three methods read deliberately different signals, so that agreement among them carries more
weight than any single method alone.

\textbf{Method 1: counting authors across the document} \cite{nath2019}. This method treats a
document as a whole. It divides the text into paragraph-sized windows, describes each window by how
often the 50 most common words occur in it, and groups similar windows together. The number of groups
is read as the number of authors, which we reduce to a binary decision: one author, or more than one.
Its check is balanced accuracy across the four known cells, which must exceed the level expected by
chance. Balanced accuracy here is the average of two quantities, how often the method is correct on
the single-author cells and how often it is correct on the multi-author cells, which prevents an
uneven mix of the two classes from inflating the score.

\textbf{Method 2: detecting style changes between sentences} \cite{schmidt2025}. This method operates
sentence by sentence. For each pair of neighbouring sentences it estimates the probability that the
two were written by different authors, using an encoder trained to capture writing style while
ignoring topic \cite{patel2025}. Longer documents contain more sentence pairs, which affects the raw
score, so we adjust for length before comparing cells. Each document receives a single number, which
we call its \emph{lean}, indicating whether its style-change pattern lies closer to the dialogue cells
or to the monologue cells. Its check is that the dialogue cells must lean more toward dialogue than
the monologue cells, by a clear margin (Cliff's $\delta$ of at least $0.33$; Cliff's $\delta$ is a
rank-based measure of how cleanly two groups separate, with a size of $0$ when they overlap completely
and $1$ when they do not overlap at all).

\textbf{Method 3: measuring question-and-answer structure} \cite{duran2022,schegloff1973}. This
method reads the text as a sequence of conversational moves. A language model tags each utterance with
its function using CAMS, a standard scheme whose labels include a question, an answer, an agreement,
or a suggestion. From these tags we measure how often an \emph{adjacency pair} is completed. An
adjacency pair is a two-part exchange in which the first part makes the second one due, such as a
question followed by an answer, or a proposal followed by its acceptance. A single writer has little
reason to produce completed pairs, so their presence is evidence of genuine dialogue. Method~3 must
clear two checks, and is the only method to which this applies. The first is the same check the others
face: its completion rate must separate the dialogue cells from the monologue cells. The second
applies to Method~3 alone, because its labels are produced by a language model acting as a tagger. For
that check we compare the tagger's labels with a human annotator's and require sufficient agreement.

\section{Results}
\textbf{Method 1 passes its check, but by a narrow margin.} Its balanced accuracy is $0.6088$ (95\%
confidence interval $0.596$ to $0.621$), above the chance level of $0.50$. The composition of this
score is important. On single-author text the method is nearly perfect, correctly labelling $0.975$
of the human monologue documents and $0.998$ of the AI monologue documents. On multi-author text it
is far weaker, recovering $0.423$ of the AI debates and none of the human debates. In sum, the method
confirms a single author reliably but detects multiple authors poorly. A monologue-side placement
therefore indicates the absence of a detected authorship signal, and should not be taken as positive
evidence of a single voice. This caveat applies to its verdict on the traces.

\textbf{Method 2 passes cleanly.} The dialogue cells lean more toward dialogue than the monologue
cells (average lean $0.0125$ versus $-0.0105$), and the separation is strong, with Cliff's
$\delta=0.6211$, well above the $0.33$ floor. The traces themselves lean further toward dialogue than
either AI anchor.

\textbf{Method 3 requires closer examination, since it must clear two checks.} On the first check,
which bears directly on our question, it performs well. Its completion rate separates dialogue from
monologue in the expected direction (Cliff's $\delta=0.4717$, above the $0.33$ floor). Within each
producer, the dialogue cell completes more pairs than the monologue cell: $0.235$ against $0.096$ for
the human cells, and $0.025$ against $0.003$ for the AI cells. The tagger therefore performs the
essential task, distinguishing genuine dialogue from genuine monologue.

On the second check, the tagger falls short of the threshold set in advance. We compared its labels
with a human annotator's and measured agreement using Cohen's $\kappa$, an agreement score that
corrects for the rate at which two labellers would match by chance. The tagger reached $\kappa=0.54$,
below the required $0.70$.

We report Method~3's result despite this outcome, and explain the decision here, since it departs from
our original plan. Two considerations motivate it. The first is the separation just described: the
check that bears directly on our question, distinguishing monologue from dialogue, is the one the
tagger passed, and passed clearly. The second is that the $0.70$ threshold is too strict for this
labelling scheme. CAMS uses a set of 28 fine-grained labels, many of them close relatives, such as
\texttt{inform} and \texttt{suggest}, which are easily confused even by careful human coders.
Unweighted $\kappa$ penalises a near miss between two similar tags as heavily as a disagreement
between two unrelated ones. The CAMS authors make the same point and recommend an agreement measure
that gives partial credit for near misses. Under such a measure, agreement rises to a tree-weighted
value of $0.62$, and the raw share of labels on which tagger and human agree is $0.74$. These values
remain below our original threshold, so we do not claim the labels are error-free. They indicate that
the labels are reasonable, and that the pre-registered threshold was too severe for an inherently
ambiguous scheme.

To make this dependence explicit, we state one assumption and carry it through the paper: wherever a
result relies on Method~3, it also relies on the assumption that the tagger is reliable. We state this
at each point, shade Method~3 in gray in every figure, and mark it with a dagger in every table.
Readers who accept the assumption may weigh Method~3 alongside the other two; readers who do not may
set it aside, since the main result rests on Methods~1 and~2 alone. Under the assumption, Method~3's
completion rate places the traces on the dialogue side: the trace rate of $0.035$ lies closer to the
AI dialogue anchor ($0.025$) than to the AI monologue anchor ($0.003$).

\begin{figure*}[t]
\centering
\includegraphics[width=0.90\textwidth]{fin_fig_split.pdf}
\caption{\textbf{Placement of the traces by each method} on the monologue--dialogue scale (poles at
$0$ and $1$). The shaded band is each method's margin, within which no lean is reported. Method~1 sits
at the edge of a wide band, just onto the monologue side, a marginal call. Method~2 sits well past the
dialogue pole. Method~3 also sits past the dialogue pole, and is shown in gray with a dagger to
indicate that it depends on the valid-tagger assumption of Section~4.}
\label{fig:split}
\end{figure*}

\textbf{Across the three methods} (Table~\ref{tab:verdict} and Figure~\ref{fig:split}), two of the
signals place the traces on the dialogue side and one places them just onto the monologue side.
Method~2, reading style changes between sentences, places them firmly past the dialogue anchor.
Method~3, reading completed question-and-answer pairs, places them past the dialogue anchor as well,
under the assumption above. Method~1, which reads only the coarse, whole-document sense of authorship,
is the method that leans monologue, and it does so only marginally, at the edge of its no-decision
band. The two methods that read fine-grained, turn-by-turn structure agree; the dissenting method is
the one that reads the least detail.

\begin{table}[t]
\centering
\small
\setlength{\tabcolsep}{4.5pt}
\caption{Placement of the reasoning traces by each method. The \emph{mono/dia} column is the validity
check every method must pass, separating monologue from dialogue on the known corpora. The
\emph{labels} column is a further check that applies only to Method~3, on the reliability of its
automatic tags; Method~3 did not pass it. Method~3's verdict is therefore conditional ($\dagger$) and
shown in gray.}
\label{tab:verdict}
\begin{tabular}{@{}lccl@{}}
\toprule
Method & mono/dia & labels & Verdict \\
\midrule
1: author count & pass & n/a & \textcolor{mono}{monologue (marginal)} \\
2: style switch & pass & n/a & \textcolor{dia}{dialogue} \\
\textcolor{m3gray}{3: adjacency pairs} & \textcolor{m3gray}{pass} & \textcolor{m3gray}{fail} &
\textcolor{m3gray}{dialogue $\dagger$} \\
\bottomrule
\end{tabular}
\end{table}

\section{What it means}
The claim here concerns the surface of the text, that is, the patterns a reader or a tool can observe
in the writing. It says nothing directly about what occurs inside the model. On that surface, the
picture leans toward dialogue. The traces change writing style from sentence to sentence in the manner
of dialogue, and, under the valid-tagger assumption, they complete question-and-answer pairs at a rate
above the dialogue anchor. The one method that points the other way reads the coarsest signal and
makes only a marginal call; it is a conservative detector, so its monologue reading indicates the
absence of an authorship signal, and should not be read as positive evidence of a single voice.

The scope of this claim is limited. A dialogue-leaning reading does not show that a model contains
several minds, just as a novel full of dialogue does not have several authors. It shows that the trace
is organised like a conversation, which is a measurable fact in itself, and a natural surface
counterpart to the internal multi-agent behaviour that Kim et al.~\cite{kim2026} describe from the
inside, through interpretability. The outside view and the inside view point the same way, which is
reassuring, though the present evidence stands on its own as a statement about surface structure.

The study supports two readings, depending on how far the tagger is trusted. Under the valid-tagger
assumption, two independent structural signals agree that the traces are dialogue-like, and the case
is fairly strong. Without that assumption, the evidence rests on Methods~1 and~2, which disagree: one
marginal monologue call from a coarse detector, and one clear dialogue call from a finer one. Even in
the more cautious reading, nothing points cleanly to monologue, and the sharpest single signal points
to dialogue. We present both readings so that they can be weighed separately.

Two caveats temper the dialogue placements, and both concern how a distance was measured rather than
its direction. Method~2's distance to the monologue side is partly an extrapolation, because the
traces are much longer than the monologue anchor, so it is more accurate to say that they track the
dialogue anchor closely over the range where the two overlap. Method~3's verdict rests on its
completion signal, since its coarser tag-frequency profiles cannot separate the two AI anchors.
Neither caveat changes a direction; each limits only how far past the pole the placement should be
read.

\section{Limitations}
Several limitations should be noted. The target is a single reasoning model, DeepSeek-R1, on a single
family of puzzles, so the verdict concerns this setting and does not yet generalise to reasoning
traces at large. The two anchors were generated with a different base model (Qwen-2.5-7B) than the
target, so part of the measured distance may reflect the change of model as well as voice structure,
and the placement is best read as ``closer to which pole'' than as an exact distance. Method~2's
monologue comparison rests on extrapolation, as noted above. Method~3's labels come from an automatic
tagger, which is why we state the valid-tagger assumption explicitly and mark its results throughout.
The validity check is the basis for our confidence: a method that clears it has shown that it can
distinguish the two structures on data whose labels are known.

\begin{thebibliography}{99}
\small
\setlength{\itemsep}{0.6pt}\setlength{\parskip}{0pt}
\bibitem{guo2025} D.\ Guo et al. DeepSeek-R1 incentivizes reasoning in LLMs through reinforcement
learning. \textit{Nature} 645:633--638, 2025.
\bibitem{kambhampati2025} S.\ Kambhampati et al. Stop anthropomorphizing intermediate tokens as
reasoning/thinking traces! \textit{arXiv:2504.09762}, 2025.
\bibitem{kim2026} J.\ Kim, S.\ Lai, N.\ Scherrer, B.\ Ag\"uera y Arcas, J.\ Evans. Reasoning models
generate societies of thought. \textit{arXiv:2601.10825}, 2026.
\bibitem{mercier2017} H.\ Mercier, D.\ Sperber. \textit{The Enigma of Reason.} Harvard Univ.\
Press, 2017.
\bibitem{minsky1986} M.\ Minsky. \textit{The Society of Mind.} Simon \& Schuster, 1986.
\bibitem{wurgaft2025} D.\ Wurgaft et al. Scaling up the think-aloud method. \textit{Proc.\ CogSci},
2025.
\bibitem{karadzhov2023} G.\ Karadzhov, T.\ Stafford, A.\ Vlachos. DeliData: a dataset for
deliberation in multi-party problem solving. \textit{Proc.\ ACM HCI} 7(CSCW2), 2023.
\bibitem{yao2023} S.\ Yao et al. Tree of thoughts: deliberate problem solving with LLMs.
\textit{NeurIPS} 36, 2023.
\bibitem{nath2019} S.\ Nath. Style change detection by threshold-based and window-merge clustering.
\textit{CLEF Working Notes}, 2019.
\bibitem{schmidt2025} G.\ Schmidt et al. Style change detection using a sequential sentence-pair
classifier. \textit{CLEF Working Notes}, 2025. \textit{arXiv:2508.00675}.
\bibitem{patel2025} A.\ Patel et al. StyleDistance: content-independent style embeddings with
synthetic parallel examples. \textit{NAACL}, 2025.
\bibitem{duran2022} N.\ Duran, S.\ Battle, J.\ Smith. Inter-annotator agreement using the
Conversation Analysis Modelling Schema. \textit{Comm.\ Methods \& Measures} 16(3):182--214, 2022.
\bibitem{schegloff1973} E.\ Schegloff, H.\ Sacks. Opening up closings. \textit{Semiotica}
8(4):289--327, 1973.
\end{thebibliography}

\end{document}
